\section{Proposed Offline Planner}
\label{sec:algorithm}

\subsection{Planner Decomposition}

The proposed planner decomposes the full problem into four stages:
\begin{enumerate}
  \item target sampling on the station model;
  \item safe candidate viewpoint generation;
  \item coverage-aware viewpoint selection and ordering;
  \item CW-feasible trajectory generation between selected viewpoints.
\end{enumerate}

This decomposition is chosen because it produces interpretable paper metrics:
coverage ratio, uncovered target set, planned $\Delta v$, executed $\Delta v$,
minimum clearance, planning time, and mission duration.

\subsection{Target Sampling}

For mesh-based experiments, the station surface is sampled into target points
$p_i$ with normals $n_i$. The sampling density should be configurable by a
target spacing parameter. For fair comparison, the same sampled target set is
used by all methods in a scenario.

\subsection{Candidate Viewpoint Generation}

A candidate viewpoint $z_j$ is generated on offset shells around the station:
\begin{equation}
  z_j = p_i + \rho_j n_i,
\end{equation}
where $\rho_j\in[\rho_{\min},\rho_{\max}]$ and the candidate is accepted only
if it satisfies keep-out constraints and does not lie inside the station
safety envelope. Candidate viewpoints can also be generated from cylindrical
or spherical shells surrounding major structural components.

For each accepted candidate $z_j$, a binary visibility vector
$A_{ij}\in\{0,1\}$ is computed:
\begin{equation}
  A_{ij}=1 \quad \Longleftrightarrow \quad
  \text{target } i \text{ is visible from candidate } j.
\end{equation}

\subsection{Viewpoint Selection}

The viewpoint subset can be selected by a weighted set-cover objective:
\begin{align}
  \min_{\alpha_j\in\{0,1\}}
  \quad &
  \sum_j c_j \alpha_j \\
  \text{s.t.}\quad &
  \sum_j A_{ij}\alpha_j \ge 1,
  \quad \forall i \in \mathcal{T}_{\mathrm{req}},
\end{align}
where $c_j$ penalizes fuel estimate, poor safety margin, and view quality. If
full set cover is too large, a greedy submodular approximation is used:
\begin{equation}
  j^\star =
  \arg\max_j
  \left[
    w_{\mathrm{cov}}\Delta C_j
    - w_d d(r,z_j)
    - w_u \hat{\Delta v}(r,z_j)
    + w_m m_j
  \right],
\end{equation}
where $\Delta C_j$ is new coverage, $d(r,z_j)$ is travel distance,
$\hat{\Delta v}$ is an estimated transfer cost, and $m_j$ is safety margin.

\subsection{Ordering and CW Transfer Generation}

After selecting candidate viewpoints, their order is computed by solving a
traveling-salesman-like problem with edge costs estimated from CW transfers.
For each edge $(i,j)$, the planner computes a dynamically feasible transfer
under bounded acceleration and rejects edges that violate safety margins.

\begin{algorithm}[t]
\caption{SOOA-based orbital coverage planner}
\label{alg:offline-planner}
\begin{algorithmic}[1]
\Require Structure geometry $\mathcal{S}$, initial state $x_0$, camera model,
HCW model, safety limits, $C_{\mathrm{req}}$, $C_{\mathrm{stop}}$
\Ensure Planned trajectory $\mathcal{P}$, attitude stream, SOOA sequence, and
inspection metrics
\State Sample area-weighted surface targets
$\mathcal{T}=\{(p_i,n_i,A_i)\}_{i=1}^N$
\State Generate viewpoint and attitude seeds around the structure
\State For each seed, generate a finite-duration terminal/dwell candidate arc
under the HCW model
\State Reject arcs violating terminal error, input, speed, clearance, or
passive-audit constraints
\State Evaluate range, FOV, incidence, dwell, and LOS visibility for accepted
arcs
\State Construct the SOOA library $\mathcal{A}$ with visibility sets
$\Gamma(a_j)$, costs $\Delta v_j$, clearances $d^{\min}_j$, and margins
$\rho_j$
\State Initialize covered set $\mathcal{C}\gets\emptyset$, selected sequence
$\sigma\gets\emptyset$, and predicted state $x\gets x_0$
\While{$C(\mathcal{C})<C_{\mathrm{stop}}$ and the arc budget is not exhausted}
  \State Compute new weighted coverage gain $\Delta C_j$ for feasible outgoing
  SOOAs
  \State Score each arc using coverage gain, transfer cost, control effort, and
  safety margin
  \State Select $a^\star=\arg\max_j s_j$
  \State Append $a^\star$ to $\sigma$ and update the predicted terminal state
  \State Update covered set $\mathcal{C}$ using $\Gamma(a^\star)$
\EndWhile
\State Roll out and audit the selected SOOA sequence under HCW dynamics
\State Export trajectory, control, attitude, FOV, coverage, safety, and planner
logs
\State \Return $\mathcal{P}$, attitude stream, SOOA sequence, and metrics
\end{algorithmic}
\end{algorithm}

\begin{algorithm}[t]
\caption{Certified reduced-SOOA-graph dynamic program}
\label{alg:certified-dp}
\begin{algorithmic}[1]
\Require SOOA graph $\mathcal G$, visibility sets $\Gamma(a_j)$, edge costs
$\ell_{jk}$, feasibility flags $\eta_{jk}$, $C_{\mathrm{req}}$, $L_{\max}$
\Ensure Minimum-cost feasible sequence on the reduced graph, or no certificate
\State Precompute target masks for all SOOAs
\State Initialize $D(\emptyset,s)\gets0$ at the initial spacecraft state $s$
\For{$\ell=0,\ldots,L_{\max}-1$}
  \ForAll{reachable states $(B,j)$ with $|B|=\ell$}
    \ForAll{$k\notin B$ with $\eta_{jk}=1$ and new visible targets}
      \State $B^+\gets B\cup\{k\}$
      \State $D(B^+,k)\gets\min\{D(B^+,k),D(B,j)+\ell_{jk}\}$
      \State Store predecessor if the value improves
    \EndFor
  \EndFor
\EndFor
\State Select the least-cost state whose weighted coverage exceeds
$C_{\mathrm{req}}$
\State Reconstruct the predecessor chain and report certificate status
\end{algorithmic}
\end{algorithm}


\subsection{Trajectory Output}

The offline planner outputs a timestamped reference trajectory
\begin{equation}
  \mathcal{P} =
  \{(t_k,r_k^{\mathrm{ref}},v_k^{\mathrm{ref}},u_k^{\mathrm{ref}})\}_{k=0}^{K}.
\end{equation}
This trajectory is saved before ROS execution and then replayed through the
controller/dynamics stack. This is the proper order for the final paper:
trajectory first, Gazebo validation second.
